\documentclass[11pt,a4paper]{article}
% Preâmbulo comum. Não duplicar pacotes nos módulos.
\usepackage[utf8]{inputenc}
\usepackage[T1]{fontenc}
\usepackage[brazil]{babel}
\usepackage{amsmath,amssymb,amsthm,mathtools}
\usepackage{geometry}
\geometry{a4paper,margin=2.3cm}
\usepackage{microtype}
\usepackage{enumitem}
\usepackage{graphicx}
\usepackage{xcolor}
\usepackage{booktabs}
\usepackage{hyperref}
\hypersetup{colorlinks=true,linkcolor=blue!40!black,urlcolor=blue!50!black}

\setcounter{tocdepth}{2}
\numberwithin{equation}{section}

\theoremstyle{plain}
\newtheorem{teorema}{Teorema}[section]
\newtheorem{proposicao}[teorema]{Proposição}
\newtheorem{lema}[teorema]{Lema}
\newtheorem{corolario}[teorema]{Corolário}

\theoremstyle{definition}
\newtheorem{definicao}[teorema]{Definição}
\newtheorem{exemplo}[teorema]{Exemplo}
\newtheorem{exercicio}{Exercício}[section]

\theoremstyle{remark}
\newtheorem{observacao}[teorema]{Observação}
\newtheorem{contraexemplo}[teorema]{Contraexemplo}
\newtheorem{intuicao}[teorema]{Intuição}

\newcommand{\R}{\mathbb{R}}
\newcommand{\N}{\mathbb{N}}
\newcommand{\Z}{\mathbb{Z}}
\newcommand{\Q}{\mathbb{Q}}
\DeclareRobustCommand{\eps}{\varepsilon}
\DeclareRobustCommand{\vet}[1]{\vec{#1}}
\DeclareRobustCommand{\dif}{\mathrm{d}}
\DeclareRobustCommand{\abs}[1]{\left\lvert #1 \right\rvert}
\DeclareRobustCommand{\norm}[1]{\left\lVert #1 \right\rVert}
\DeclareMathOperator{\grad}{grad}
\DeclareMathOperator{\Div}{div}
\DeclareMathOperator{\rot}{rot}

\newcommand{\base}[1]{\par\medskip\noindent\textbf{Base para relembrar.} #1\par\medskip}
\newcommand{\cuidado}[1]{\par\medskip\noindent\textbf{Quando a hipótese falha.} #1\par\medskip}

\newcommand{\figura}[2]{%
  \begin{figure}[htbp]
  \centering
  \IfFileExists{#1}{%
    \includegraphics[width=0.84\textwidth]{#1}%
  }{%
    \fbox{\parbox{0.8\textwidth}{\centering Figura ainda não gerada:\\ \texttt{#1}}}%
  }
  \caption{#2}
  \end{figure}%
}

\newenvironment{esqueleto}{%
  \par\medskip\noindent\textbf{Esqueleto da demonstração.}\par\nopagebreak
  \begin{enumerate}[leftmargin=1.4em]
}{\end{enumerate}\medskip}

\title{Gabarito: otimização e multiplicadores de Lagrange}
\author{}
\date{}
\begin{document}
\maketitle
\tableofcontents

\section{Exercícios}

\begin{exercicio}[{[Básico --- ponto crítico e teste]}]
Classificar $x^2+y^2-4x+2y+5$.

\textbf{Solução.}
$f_x=2x-4$ e $f_y=2y+2$. O sistema $x=2$, $y=-1$ tem solução única $(2,-1)$. As segundas parciais são constantes, $f_{xx}=2$, $f_{yy}=2$, $f_{xy}=0$, então $D=4>0$ e $f_{xx}>0$. O teste dá mínimo local estrito. O valor é
\[
f(2,-1)=4+1-8-2+5=0.
\]
Completando quadrados, $f(x,y)=(x-2)^2+(y+1)^2$, que é $\ge 0$ em todo o plano. O mínimo local é global, e não há máximo.
\end{exercicio}

\begin{exercicio}[{[Básico --- sela e mínimo local]}]
Classificar $x^3+y^3-3xy$.

\textbf{Solução.}
$f_x=3x^2-3y$ e $f_y=3y^2-3x$. Então $y=x^2$ e $x=y^2$. Substituindo, $x=(x^2)^2=x^4$, logo $x^4-x=x(x^3-1)=0$. As raízes reais são $x=0$ e $x=1$, porque $x^3=1$ só tem a raiz real $1$. Os pontos críticos são $(0,0)$ e $(1,1)$.

As segundas parciais são $f_{xx}=6x$, $f_{yy}=6y$ e $f_{xy}=-3$. Assim
\[
D=36xy-9.
\]
Em $(0,0)$, $D=-9<0$: sela. Em $(1,1)$, $D=27>0$ e $f_{xx}=6>0$: mínimo local estrito, com
\[
f(1,1)=1+1-3=-1.
\]
Esse mínimo não é global. Por exemplo $f(-2,0)=-8<-1$, e $f(x,0)=x^3\to-\infty$ quando $x\to-\infty$.
\end{exercicio}

\begin{exercicio}[{[Básico --- hessiana]}]
Natureza de $x^2-4xy+y^2$ na origem.

\textbf{Solução.}
$f_x=2x-4y$ e $f_y=-4x+2y$. O sistema linear tem matriz de coeficientes com determinante $4-16=-12\neq 0$, então a única solução é $(0,0)$. A hessiana é constante,
\[
H=\begin{pmatrix}2&-4\\-4&2\end{pmatrix},\qquad D=4-16=-12<0.
\]
O teste dá sela. Ao longo de $y=0$, $f(x,0)=x^2$, positivo para $x\neq 0$. Ao longo de $y=x$, $f(x,x)=x^2-4x^2+x^2=-2x^2$, negativo para $x\neq 0$. Os dois sinais numa vizinhança qualquer da origem confirmam a sela sem depender só do nome do teste.
\end{exercicio}

\begin{exercicio}[{[Básico --- uma restrição afim]}]
Extremos de $x^2+y^2$ na reta $x+y=1$.

\textbf{Solução.}
Ponhamos $g(x,y)=x+y$. Então $\nabla g=(1,1)\neq\vet{0}$ em todo o plano, em particular na reta. O sistema é $(2x,2y)=\lambda(1,1)$ e $x+y=1$. Logo $2x=2y$, então $x=y$, e $x=y=1/2$. O valor é
\[
f\Bigl(\frac12,\frac12\Bigr)=\frac14+\frac14=\frac12.
\]
Não há outro candidato. Sobre a reta, $y=1-x$ e
\[
h(x)=x^2+(1-x)^2=2x^2-2x+1.
\]
Quando $\abs{x}\to\infty$, $h(x)\to\infty$. O conjunto $\{(x,1-x):h(x)\le 1\}$ é fechado e limitado, logo o mínimo existe e tem de ser o único ponto crítico da restrição. O mínimo é $1/2$ em $(1/2,1/2)$. Não há máximo: $h(x)$ não é limitado superiormente. Weierstrass não se aplica à reta inteira porque ela não é limitada.
\end{exercicio}

\begin{exercicio}[{[Intermediário --- compacto com interior]}]
Extremos de $x^2+y^2+x$ no quadrado $[-1,1]^2$.

\textbf{Solução.}
O quadrado é fechado e limitado, e $f$ é contínua, então máximo e mínimo absolutos existem. No interior, $f_x=2x+1$ e $f_y=2y$. O único ponto crítico é $(-1/2,\,0)$, que está no interior do quadrado. O valor é
\[
f\Bigl(-\frac12,0\Bigr)=\frac14-\frac12=-\frac14.
\]
A hessiana tem $D=4$ e $f_{xx}=2>0$, então esse ponto é mínimo local. Para o absoluto, comparamos com a fronteira.

No lado $x=-1$, $-1\le y\le 1$: $f=1+y^2-1=y^2$, que vai de $0$ a $1$.
No lado $x=1$: $f=1+y^2+1=2+y^2$, que vai de $2$ a $3$.
No lado $y=1$: $g(x)=x^2+x+1$ em $[-1,1]$. Então $g'(x)=2x+1=0$ em $x=-1/2$, e $g(-1/2)=1/4-1/2+1=3/4$. Nos cantos, $g(-1)=1$ e $g(1)=3$.
No lado $y=-1$ a função de $x$ é a mesma, $x^2+x+1$, com os mesmos valores $3/4$, $1$ e $3$.

Os números que aparecem são $-1/4$, $0$, $3/4$, $1$, $2$ e $3$. O mínimo absoluto é $-1/4$, em $(-1/2,\,0)$. O máximo absoluto é $3$, em $(1,1)$ e em $(1,-1)$.
\end{exercicio}

\begin{exercicio}[{[Intermediário --- circunferência]}]
Extremos de $x+2y$ em $x^2+y^2=4$.

\textbf{Solução.}
A circunferência é compacta e $f$ é contínua, então máximo e mínimo existem. Com $g=x^2+y^2$, tem-se $\nabla g=(2x,2y)$, que se anula só na origem, fora da circunferência. O sistema é
\[
(1,2)=\lambda(2x,2y),\qquad x^2+y^2=4.
\]
Logo $x=1/(2\lambda)$ e $y=1/\lambda$, então $y=2x$. Substituindo, $x^2+4x^2=4$, $5x^2=4$, $x=\pm 2/\sqrt{5}$ e $y=\pm 4/\sqrt{5}$, com os sinais iguais. Os valores são
\[
f=x+2y=x+4x=5x=\pm\frac{10}{\sqrt{5}}=\pm 2\sqrt{5}.
\]
O máximo é $2\sqrt{5}$, em $\bigl(2/\sqrt{5},\,4/\sqrt{5}\bigr)$. O mínimo é $-2\sqrt{5}$, no ponto oposto.

Os pontos também se escrevem $\bigl(\pm 2\sqrt{5}/5,\,\pm 4\sqrt{5}/5\bigr)$.
\end{exercicio}

\begin{exercicio}[{[Intermediário --- teste inconclusivo]}]
Pontos críticos de $x^2 y$ e a natureza de três deles.

\textbf{Solução.}
\begin{enumerate}[label=\alph*)]
\item $f_x=2xy$ e $f_y=x^2$. A segunda equação força $x=0$, e a primeira fica automática. O conjunto crítico é o eixo $y$, isto é, todos os pontos $(0,y)$. As segundas parciais são $f_{xx}=2y$, $f_{yy}=0$ e $f_{xy}=2x$. Em $x=0$, $D=0\cdot 0-(0)^2=0$. O teste não classifica nenhum desses pontos.

\item Perto de $(0,0)$, $f(0{,}1,\,0{,}1)=(0{,}1)^2(0{,}1)>0$ e $f(0{,}1,\,-0{,}1)<0$. A origem não é extremo local.

Perto de $(0,1)$, escreva $x=h$ e $y=1+k$. Então $f=h^2(1+k)$. Se $\abs{k}<1/2$, tem-se $1+k>0$, logo $f\ge 0=f(0,1)$. O ponto é mínimo local, não estrito ao longo do eixo, onde $f$ continua zero.

Perto de $(0,-1)$, $f=h^2(-1+k)$. Se $\abs{k}<1/2$, então $-1+k<0$ e $f\le 0=f(0,-1)$. O ponto é máximo local.
\end{enumerate}
O mesmo número $D=0$ cobre um ponto que não é extremo, um mínimo local e um máximo local. A hessiana não separa os três casos; o sinal de $y$ na fatoração $x^2 y$ separa.
\end{exercicio}

\begin{exercicio}[{[Avançado --- direções da sela]}]
Sela de $x^2+4xy+y^2$ com vetores explícitos.

\textbf{Solução.}
\begin{enumerate}[label=\alph*)]
\item $f_x=2x+4y$ e $f_y=4x+2y$. A matriz do sistema é $\begin{pmatrix}2&4\\4&2\end{pmatrix}$, de determinante $-12\neq 0$. A única solução é $(0,0)$. As segundas parciais são $f_{xx}=2$, $f_{yy}=2$ e $f_{xy}=4$, então
\[
D(0,0)=4-16=-12<0.
\]
A forma quadrática da hessiana é $Q(h,k)=2h^2+8hk+2k^2$, e $f=Q/2$. As segundas parciais são constantes, então no desenvolvimento de Taylor o resto além de $Q$ é zero: $f(h,k)-f(0,0)=\frac12 Q(h,k)$.

\item Para $\vet{u}=(1,0)$, que é unitário, $Q(\vet{u})=2>0$ e $f(t,0)=t^2>0$ para $t\neq 0$.
Para a direção de $(1,-1)$, o vetor unitário é $\vet{v}=(1/\sqrt{2},\,-1/\sqrt{2})$. Primeiro $Q(1,-1)=2-8+2=-4$, logo
\[
Q(\vet{v})=\frac{1}{2}Q(1,-1)=-2<0.
\]
Ao longo dessa reta, $f(t,-t)=t^2-4t^2+t^2=-2t^2<0$ para $t\neq 0$. Em toda bola há valores positivos e valores negativos. A origem é sela. O sinal de $D$ anunciou o fato; os dois vetores são a prova direta, porque o resto é zero e o sinal de $f$ é o sinal de $Q$.
\end{enumerate}
\end{exercicio}

\begin{exercicio}[{[Avançado --- duas restrições]}]
Ponto da reta $x+y+z=1$, $x+y-z=0$ mais próximo da origem.

\textbf{Solução.}
Os gradientes $\nabla g=(1,1,1)$ e $\nabla h=(1,1,-1)$ são linearmente independentes em todo o $\R^3$: se $a(1,1,1)+b(1,1,-1)=\vet{0}$, então $a+b=0$ e $a-b=0$, logo $a=b=0$. A hipótese do método vale em toda a reta.

O sistema $\nabla f=\lambda\nabla g+\mu\nabla h$ é
\begin{align*}
2x&=\lambda+\mu,\\
2y&=\lambda+\mu,\\
2z&=\lambda-\mu,
\end{align*}
com $x+y+z=1$ e $x+y-z=0$. As duas primeiras equações dão $x=y$. Somando as restrições, $2(x+y)=1$; subtraindo, $2z=1$. Logo $z=1/2$ e $x+y=1/2$. Com $x=y$, vem $x=y=1/4$. O ponto é $(1/4,\,1/4,\,1/2)$. Os multiplicadores saem de $\lambda+\mu=1/2$ e $\lambda-\mu=1$: $\lambda=3/4$ e $\mu=-1/4$. O quadrado da distância é
\[
f=\frac{1}{16}+\frac{1}{16}+\frac14=\frac{3}{8}.
\]
A distância é $\sqrt{3/8}=\sqrt{6}/4$.

Conferência por parametrização. Um vetor diretor é o produto vetorial
\[
(1,1,1)\times(1,1,-1)=(-2,2,0),
\]
paralelo a $(-1,1,0)$. A reta é
\[
\Bigl(\frac14-t,\,\frac14+t,\,\frac12\Bigr),\qquad t\in\R.
\]
O quadrado da distância fica
\[
\Bigl(\frac14-t\Bigr)^2+\Bigl(\frac14+t\Bigr)^2+\Bigl(\frac12\Bigr)^2
=2\Bigl(\frac{1}{16}+t^2\Bigr)+\frac14
=\frac{3}{8}+2t^2.
\]
O mínimo é $3/8$, em $t=0$, e não há outro. Quando $\abs{t}\to\infty$ a expressão tende a infinito, então não há máximo. A reta é fechada e não é limitada; Weierstrass não garante extremo, e de fato o máximo não existe. O mínimo existe por causa do comportamento no infinito, não por compacidade da reta.
\end{exercicio}

\end{document}
